\documentclass[11pt]{amsart}
\usepackage[margin=1.1in]{geometry}
\usepackage[T1]{fontenc}
\usepackage{amsmath,amssymb,amsthm,mathtools,booktabs,hyperref}
\hypersetup{colorlinks=true,linkcolor=blue!50!black,citecolor=blue!50!black,urlcolor=blue!50!black}
\usepackage[dvipsnames]{xcolor}
\numberwithin{equation}{section}
\emergencystretch=4em
\newtheorem{theorem}{Theorem}[section]
\newtheorem{lemma}[theorem]{Lemma}
\newtheorem{proposition}[theorem]{Proposition}
\newtheorem{corollary}[theorem]{Corollary}
\theoremstyle{definition}
\newtheorem{definition}[theorem]{Definition}
\theoremstyle{remark}
\newtheorem{remark}[theorem]{Remark}
\newcommand{\li}{\operatorname{li}}
\newcommand{\Rp}{\operatorname{Re}}
\newcommand{\Ip}{\operatorname{Im}}
\newcommand{\eps}{\varepsilon}
\newcommand{\code}[1]{\texttt{#1}}
\DeclareMathOperator{\sgn}{sgn}

\title[Ramanujan's prime-counting inequality]{Ramanujan's prime-counting inequality holds for all $x \geq 38\,358\,837\,683$}
\author{Jaideep Sai Padhi}
\address{Purdue University, West Lafayette, IN, USA}
\email{jpadhi@purdue.edu}
\date{\today}
\subjclass[2020]{Primary 11N05; Secondary 11M26, 11Y35}
\keywords{Ramanujan's inequality, prime-counting function, zero-density estimates, zero-free regions, explicit formula, interval arithmetic}

\begin{document}
\begin{abstract}
Ramanujan observed that $\pi(x)^2 < \frac{e x}{\log x}\,\pi\!\left(\frac{x}{e}\right)$ for all sufficiently large $x$.
Dudek and Platt showed, assuming the Riemann hypothesis, that the last counterexample is $x = 38\,358\,837\,682$.
Unconditionally the inequality was known for $38\,358\,837\,683 \le x \le e^{103}$ (Johnston) and for $x \ge e^{3158.442}$ (Axler).
We close the gap: the inequality holds unconditionally for every $x \ge 38\,358\,837\,683$.
The proof uses no verification of the Riemann hypothesis beyond the height $3\cdot 10^{12}$ established by Platt and Trudgian.
It combines a smoothed explicit formula of B\"uthe's type with three new explicit inputs:
\begin{itemize}
\item a log-free zero-density estimate near the $1$-line, obtained by reworking Bellotti's argument with a sharp asymptotic evaluation of the Barban--Vehov mean square;
\item a height-dependent zero-free region, obtained by evaluating the inequality of Bellotti, Trudgian and Yang exactly;
\item a repulsion lemma showing that two hypothetical zeros close to the $1$-line cannot lie close to each other in height.
\end{itemize}
Every numerical step is certified in rigorous interval (ball) arithmetic.
Along the way we record several corrections to the explicit inputs we use, none of which affects the published main results.
\end{abstract}
\maketitle
\tableofcontents
\section{Introduction}

In his second notebook (see \cite{BerndtIV}, p.~112) Ramanujan stated that
\begin{equation}\label{eq:ram}
\pi(x)^2 < \frac{e x}{\log x}\, \pi\!\left(\frac{x}{e}\right)
\end{equation}
holds for all sufficiently large $x$, where $\pi(x)$ is the number of primes up to $x$.
Expanding $\pi$ by the prime number theorem with the error term $O(x\log^{-6}x)$ shows that the two sides differ by
$\frac{x^2}{\log^6 x}\,(1+o(1))$, so \eqref{eq:ram} is a statement about the sixth term of the asymptotic expansion.
It requires the relative error $|\pi(x)-\li(x)|\log x/x$ to be smaller than roughly $1/(3\log^4 x)$.
Following Axler \cite{Axler}, let $N_R$ denote the smallest positive integer such that \eqref{eq:ram} holds for every real $x \ge N_R$.

\subsection*{What was known}
\begin{itemize}
\item Under the Riemann hypothesis, Dudek and Platt \cite{DudekPlatt} showed $N_R = 38\,358\,837\,683$.
  They combined Hassani's conditional bound \cite{Hassani} with a computation for $x \le 1.15\cdot 10^{16}$.
  The largest prime counterexample, $38\,358\,837\,677$, had been found by Galway (see \cite{DudekPlatt}).
\item Unconditionally, the inequality is known in two ranges:
  \begin{itemize}
  \item for $38\,358\,837\,683 \le x \le 10^{19}$ (Axler \cite{Axler2}, building on B\"uthe's bounds \cite{Buthe2018}), extended to $x\le e^{58}$ by Platt and Trudgian \cite{PlattTrudgian} and to $x \le e^{103}$ by Johnston \cite[Thm.~5.1]{Johnston};
  \item for $x \ge e^{X_0}$, where $X_0$ fell successively as follows (see \cite{Axler} for the history):
  \[
  9658 \to 9032\ \text{\cite{Axler2}} \to 3915\ \text{\cite{PlattTrudgian}} \to 3604\ \text{\cite{CullyHugillJohnston}} \to 3361\ \text{\cite{JohnstonYang}} \to 3158.442\ \text{\cite{Axler}}.
  \]
  \end{itemize}
\end{itemize}
The range $e^{103} < x < e^{3158.442}$ remained open.
Two recent preprints \cite{De1,De2} claim to settle it; we explain in \S\ref{sec:de} why their arguments are not proofs.

\subsection*{Result}
\begin{theorem}\label{thm:main}
Inequality \eqref{eq:ram} holds for every real $x$ with $e^{58} \le x \le e^{3159}$.
Consequently it holds for every $x \ge 38\,358\,837\,683$, and $N_R = 38\,358\,837\,683$ unconditionally.
\end{theorem}

The second statement follows from the first together with \cite[Thm.~5.1]{Johnston} and \cite[Thm.~1.1]{Axler}.
The value of $N_R$ also uses the computations of Galway and of Dudek--Platt \cite{DudekPlatt} for $x \le 10^{11}$.

The only information about zeros of $\zeta$ that we take from a large-scale computation is the theorem of Platt and Trudgian \cite{PlattTrudgianRH}:
every non-trivial zero $\rho$ of $\zeta$ with $0 < \Ip\rho \le H_0 := 3\cdot 10^{12}$ satisfies $\Rp \rho = 1/2$.
All remaining information about zeros comes from explicit theorems, and our own computations are certified with rigorous ball arithmetic \cite{arb}.

\subsection*{Method}
Write $L = \log x$.
\begin{enumerate}
\item By an elementary argument (\S\ref{sec:reduction}), \eqref{eq:ram} at $x$ follows from
  $|\pi(x)-\li(x)| \le \eps_R(L)\,x/L$ and $|\pi(x/e)-\li(x/e)| \le \eps_R(L)\,(x/e)/(L-1)$.
  Here $\eps_R(L) \approx 1/(3L^4)$ is an explicit function (Definition~\ref{def:epsR}).
  This in turn follows from relative bounds for $\psi(t)-t$ on $[e^{43}, x]$ (Lemma~\ref{lem:transfer}).
\item The bound for $\psi$ comes from a smoothed explicit formula with Logan's kernel, as in B\"uthe \cite{Buthe2016,Buthe2018} (\S\ref{sec:explicit}).
  The zeros with $|\gamma| \le H_0$ lie on the critical line.
  Zeros above $H_0$ are controlled by three ingredients:
  \begin{enumerate}
  \item[(a)] a zero-free region $\sigma > 1 - \lambda^*(\log t)/\log t$ (\S\ref{sec:zfr});
  \item[(b)] explicit zero-density estimates $N(\sigma,T)$ (\S\ref{sec:density});
  \item[(c)] a repulsion lemma for pairs of zeros close to the $1$-line (\S\ref{sec:rep}), which feeds into (b).
  \end{enumerate}
\item The binding range is $L \approx 900$, where the kernel has bandwidth $T_1 \approx e^{30}$.
  There the decisive zeros are those with $1-\beta \approx 0.21/\log\gamma$, at heights just above $H_0$.
  With only the published zero-free region and density constants, the zeros above $H_0$ alone would exceed the available error budget several times over in this range.
\end{enumerate}

\subsection*{New ingredients}
\begin{enumerate}
\item \textbf{Zero density} (Theorem~\ref{thm:density}).
  We rework Bellotti's log-free zero-density estimate \cite{Bellotti} for $\sigma$ near $1$ and $3\cdot 10^{12} < T \le e^{40}$.
  The main gain comes from an asymptotically sharp, fully explicit evaluation of the Barban--Vehov mean square $\sum_n \Psi(n)^2 e^{-n/X}/n$ (Theorem~\ref{thm:G}).
  This is the quantity Bellotti calls $d_4$; our bound is about $1.5$ where her Table~5 gives $85$--$105$.
  The bound uses a smoothed contour argument with a certified enclosure of $1/\zeta$ in a box.
  Further, smaller improvements:
  \begin{itemize}
  \item a corrected and sharpened bound for the off-diagonal pair sum, using a crowding argument;
  \item a log-free pointwise bound for the mollifier;
  \item a sharper treatment of the contour integral of $\zeta(1/2+it)$;
  \item parameters optimised separately for each $(\sigma,T)$-cell.
  \end{itemize}
\item \textbf{A height-dependent zero-free region} (Theorem~\ref{thm:H}).
  Bellotti, Trudgian and Yang \cite{BTY} proved $\zeta(\sigma+it)\neq 0$ for $\sigma > 1 - 1/(4.896\log t)$.
  We evaluate their basic inequality exactly at every height and iterate.
  This gives $(1-\beta)\log\gamma \ge 0.21384$ for $\log\gamma \le 29$, which is $R = 4.676$ in place of $4.896$.
  Several of their constants are tied to $\sigma \ge 0.99346$, so we re-derive all of them for the smaller $\sigma$ and the larger $\eta$ we need.
\item \textbf{Repulsion} (Theorem~\ref{thm:rep}).
  Keeping a second zero, together with its reflection, in the zero sum of the same inequality shows the following.
  Two zeros $\beta_j+i\gamma_j$ with $(1-\beta_j)\log\gamma_j \le 0.35$ at heights near $e^{29}$ satisfy $|\gamma_0-\gamma_1| \ge 0.058$.
  This removes the near-diagonal part of the pair sum in the density estimate.
\end{enumerate}

\subsection*{Certification}
The zero-free and repulsion certificates, the constants of \cite{BTY} that we re-derive, the box enclosure of $1/\zeta$, the reduction of \S\ref{sec:reduction} and the smoothing constant are certified with ball arithmetic (\code{python-flint}/Arb \cite{arb}).
The final scan over $L$, which evaluates explicit upper sums of monotone majorants, runs in double precision; \S\ref{sec:computation} explains why its rounding errors are far below its margin.
The code, tables and logs are available at \url{https://github.com/jaideepsaipadhi/ramanujans-inequality}.
Section~\ref{sec:computation} lists the certificate for each lemma.

\subsection*{Corrections to the literature}
We found the following points while re-deriving the explicit inputs we use; none of them affects the main published results.
\begin{itemize}
\item In \cite{BTY}:
  \begin{itemize}
  \item a normalisation slip in the final inequality: the Lemma-10 term should be multiplied by $\kappa = 433/859$;
  \item an incorrect value in their Table~1, at $p = 5$;
  \item a digamma majorant whose integral is larger than stated;
  \item a positivity test that is valid only for $\sigma \ge 0.99346$.
  \end{itemize}
\item In B\"uthe's smoothing estimate \cite[Prop.~4]{Buthe2016} (quoted as Prop.~1 in \cite{Buthe2018}), the asymptotic quoted for $\nu_c$ does not match the stated normalisation of the kernel; our own proof of the smoothing step (Lemma~\ref{lem:smooth}) avoids the issue.
\end{itemize}
Details are in \S\ref{sec:remarks}.
\section{Reduction to bounds for $\psi$}\label{sec:reduction}

\begin{definition}\label{def:epsR}
For $L > 1$ and $x = e^L$, put
\[
g(L) = \frac{\li(x)^2 - \frac{e x}{L}\li(x/e)}{(x/L)^2}, \qquad b(L) = \frac{2\li(x)}{x/L} + \frac{L}{L-1}.
\]
Define $\eps_R(L) := \big(-b + \sqrt{b^2 - 4g}\big)/2$ whenever $g(L) < 0$.
\end{definition}

Asymptotically $\eps_R(L) = \frac{1}{3L^4}\big(1+O(L^{-1})\big)$; for example $\eps_R(L)\cdot 3L^4$ equals $1.2706$ at $L = 58$, $1.0145$ at $L = 907$ and $1.0041$ at $L = 3159$.
Computing $\eps_R$ directly is delicate, because $g$ is the difference of two quantities of size $1$ that agree to within $L^{-4}$.

\begin{lemma}\label{lem:epsR}
Put $A(L) = L e^{-L}\li(e^L)$, $u = 1/L$ and $P(u) = \sum_{k<16} k!\,u^k$.
Then $g = A(L)^2 - \frac{L}{L-1}A(L-1)$, $b = 2A(L) + \frac{L}{L-1}$, and $\eps_R = -2g/(b + \sqrt{b^2-4g})$.
Moreover $A(L) = P(u) + R(L)$, where
\[
-6e^{-L/2} \le R(L) \le \frac{16!\,u^{16}}{1 - 34u} + 6e^{-L/2} \qquad (L \ge 57).
\]
For every cell $[a, a+0.05] \subset [58, 3160.1]$ we have $g < 0$ on the cell, and $\eps_R$ admits an explicit certified lower bound there.
\end{lemma}

\begin{proof}
The formulas for $g$ and $b$ follow by dividing the definitions by $(x/L)^2$ and $x/L$.
For $R$, integrate $\int_a^x dt/\log t$ by parts $16$ times with $a = \sqrt x$.
The boundary term at $a$ is $\li(a) - aP(2/L)/(L/2)$, which has absolute value at most $3a/(L/2)$ because $A$ and $P$ are both at most $1.1$ for arguments $\ge 28.5$.
The integral remainder satisfies $0 \le (L/x)\,16!\int_a^x dt/\log^{17}t \le 16!\,u^{16}/(1-34u)$.
To see this, note that the logarithmic derivative of $e^s s^{-17}$ is at least $1 - 34/L$ on $[L/2, L]$.

The cancelling part $S(u) = P(u)^2 - P\big(\tfrac{u}{1-u}\big)/(1-u)$ is a power series in $u$ with integer coefficients.
These are computed exactly up to order $70$, and $s_0 = \dots = s_3 = 0$, $s_4 = -1$, $s_5 = -14$.
The tail is bounded by $\big(\sum_{k<16}k!\big)(2u)^{71}/(1-2u)$.
Evaluating $S$ as $u^4\cdot(S/u^4)$ in ball arithmetic on each cell avoids the cancellation, and the remaining terms involve $R$ only.
This is carried out by \code{certify/epsR\_cert.py} at $200$ bits.
Point values agree with a $60$-digit direct evaluation; the enclosures have relative width at most $1.9\cdot 10^{-3}$, which is the variation of $\eps_R$ across a cell.
\end{proof}

\begin{lemma}\label{lem:reduce}
Let $L = \log x \ge 58$ and $\eps = \eps_R(L)$.
Assume $g(L) < 0$ and
\begin{equation}\label{eq:pibounds}
|\pi(x) - \li(x)| \le \eps\, \frac{x}{L}, \qquad \Big|\pi\Big(\frac xe\Big) - \li\Big(\frac xe\Big)\Big| \le \eps\, \frac{x/e}{L-1},
\end{equation}
with at least one of the two inequalities strict.
Then \eqref{eq:ram} holds at $x$.
\end{lemma}

\begin{proof}
Since $\li(x) > \eps x/L$, \eqref{eq:pibounds} gives $0 < \pi(x) \le \li(x) + \eps x/L$ and $\pi(x/e) \ge \li(x/e) - \eps x/(e(L-1))$.
It therefore suffices to show
$(\li(x) + \eps x/L)^2 < \frac{ex}{L}\big(\li(x/e) - \eps\frac{x}{e(L-1)}\big)$.
After dividing by $(x/L)^2$ this reads $g + b\eps + \eps^2 < 0$, and with $\eps = \eps_R(L)$ the left side equals $0$.
Strictness in \eqref{eq:pibounds} makes the inequality strict.
\end{proof}

Writing $y = x/e$, so that $\log y = L - 1$, the hypotheses \eqref{eq:pibounds} for all $x$ with $\log x \in [58, 3159]$ become the following two conditions:
\begin{equation}\label{eq:target}
\begin{aligned}
|\pi(y) - \li(y)| &< \eps_R(\log y + 1)\,\frac{y}{\log y} &&(57 \le \log y \le 3158),\\
|\pi(y) - \li(y)| &< \eps_R(\log y)\,\frac{y}{\log y} &&(58 \le \log y \le 3159).
\end{aligned}
\end{equation}
No monotonicity of $\eps_R$ is needed: both conditions are checked directly, cell by cell (Proposition~\ref{prop:transfer}).

\begin{lemma}[transfer]\label{lem:transfer}
Let $y \ge 2$.
Then
\[
\pi(y) - \li(y) = \frac{\vartheta(y) - y}{\log y} + \int_2^y \frac{\vartheta(t) - t}{t\log^2 t}\,dt + \frac{2}{\log 2} - \li(2),
\qquad \Big|\frac{2}{\log 2} - \li(2)\Big| < 1.84.
\]
Moreover $0 \le \psi(t) - \vartheta(t) \le 1.4262\sqrt t$ for $t > 0$ \cite[Thm.~13]{RS62}.
\end{lemma}

\begin{proof}
Partial summation gives $\pi(y) = \vartheta(y)/\log y + \int_2^y \vartheta(t)\,dt/(t\log^2 t)$.
Integration by parts gives $\li(y) = y/\log y + \int_2^y dt/\log^2 t + \li(2) - 2/\log 2$.
Subtracting the two identities proves the formula.
\end{proof}

\begin{proposition}\label{prop:transfer}
Suppose that for every certified interval $[L_1, L_2] \subset [43, 3159]$ of the final scan (\S\ref{sec:computation})
\begin{equation}\label{eq:psicert}
|\psi(t) - t| \le q_{[L_1,L_2]}\,\iota_{[L_1,L_2]}\,\frac{\eps_R(L_2+1)}{1 + 2/L_1}\, t + a(t) \qquad (\log t \in [L_1, L_2]).
\end{equation}
Here $q_{[L_1,L_2]}$ is the ratio computed by the scan for that interval, $\iota_{[L_1,L_2]} \le 1 + 10^{-3}$ is the correction factor of Lemma~\ref{lem:iota}, and $a(t) = 2 + e^{2\cdot 10^{-4}}(\log t + 10^{-4})\big(2.01\cdot 10^{-4}\sqrt t + \tfrac12\log\log(2t^2)\big)$ collects the additive terms of Lemma~\ref{lem:smooth} and Proposition~\ref{prop:explicit}.
Suppose also that $|\psi(t) - t| \le 0.94\sqrt t$ for $11 < t \le 10^{19}$ \cite[Thm.~2]{Buthe2018}.
Then \eqref{eq:target} holds.
\end{proposition}

\begin{proof}
Combine Lemma~\ref{lem:transfer} with \eqref{eq:psicert}, the bound $|\vartheta - t| \le |\psi - t| + 1.4262\sqrt t$, and the trivial bound $|\vartheta(t) - t| \le 11$ for $t \le 11$.
The integral over $[2, e^{43}]$ is at most
$\big((0.94 + 1.4262)\cdot 2 e^{21.5}/\log^2 2\big) + 121$.

On $[e^{43}, y]$ we use upper Riemann sums in $u = \log t$ with step $h = 0.05$.
On each cell we use the largest certified bound meeting it, and the inequality $\int_a^b e^u u^{-2}du \le (e^b - e^a)/a^2$.
For $y$ in a cell $[a, a+h]$ with $a \ge 57$, the numerator is bounded over the whole cell.
It is compared with the certified lower bounds of Lemma~\ref{lem:epsR} for $\eps_R$ on $[a+1, a+1+h]$ and on $[a, a+h]$; the latter comparison is needed only when $a \ge 58$.
The script \code{certify/transfer\_check.py} carries this out in $40$-digit arithmetic for every cell of $[57, 3159)$.
The largest value of
\[
\frac{\text{bound for } |\pi(y) - \li(y)|}{\min(\eps_R(\log y+1),\,\eps_R(\log y))\cdot y/\log y}
\]
that it finds is $0.99577$, attained at $\log y = 917$.
\end{proof}

\begin{remark}
The factor $(1+2/L_1)^{-1}$ in \eqref{eq:psicert} is what the integral term costs.
Heuristically, $\int^{y}\eps_R(\log t+1)\,dt/\log^2 t \approx \eps_R(\log y+1)\frac{y}{\log^2 y}(1 + 6/\log y)$, so the total is $\approx(1 + 1/L + 6/L^2) \le 1 + 2/L$.
Proposition~\ref{prop:transfer} does not rely on this heuristic; it verifies the inequality numerically, cell by cell.
\end{remark}

\begin{lemma}\label{lem:iota}
In the scan, $\eps$ is defined by three steps of the iteration $\eps_{n+1} = C/\kappa(\eps_n)$, where
$\kappa(\eps) = (L_2 + \lambda_s)/(L_1 - \lambda_s)$, $\lambda_s = \log(10\sqrt c/\eps)$, and $C = 0.7946^{-1}\,1.02^{-1}f\tau_{12}\sqrt c$.
The smoothing term it charges is $f\tau_{12}$, whereas Proposition~\ref{prop:master} requires $f\tau_{12}\,\kappa(\eps_3)/\kappa(\eps_2)$.
Let $\theta = (L_1+L_2)/((L_2+\lambda_{\max})(L_1-\lambda_{\max}))$ and $\kappa_0 = (L_2+\lambda_{\max})/(L_1-\lambda_{\max})$, where $\lambda_{\max} = \log(10\sqrt{110}/\eps_{\min})$ and $\eps_{\min} = 0.3\,\tau_{12}\sqrt{30}/(0.7946\cdot1.02\cdot 10)$.
Then
\[
\kappa(\eps_3)/\kappa(\eps_2) \le \iota := (1 + 10^{-8})\exp(\theta^3\log\kappa_0) .
\]
\end{lemma}

\begin{proof}
With $y_n = \log\eps_n$, we have $y_{n+1} = \log C - \log\kappa(e^{y_n})$.
Moreover $|\frac{d}{dy}\log\kappa| = (L_1+L_2)/((L_2+\lambda_s)(L_1-\lambda_s)) \le \theta$, since $\lambda_s\le\lambda_{\max}$.
Hence $|y_{n+1} - y_n| \le \theta^n|y_1 - y_0|$, with $y_0 = \log C$ and $|y_1 - y_0| = \log\kappa(\eps_0) \le \log\kappa_0$.
Therefore $|\log\kappa(\eps_3) - \log\kappa(\eps_2)| \le \theta|y_3 - y_2| \le \theta^3\log\kappa_0$.
The factor $1+10^{-8}$ absorbs the Logan factor $e^{\eps(c+1)/6}$ of Lemma~\ref{lem:logancomplex}, which is $\le 1 + 5\cdot10^{-10}$ whenever $T_1 > H_0$.
The script \code{certify/transfer\_check.py} evaluates $\iota$ for every interval and checks $\kappa_0 \le 10$ and $L_1 - \lambda_{\max} > 1$.
At $L_1 = 43$ it gives $\iota \le 1 + 10^{-3}$, and $\iota - 1$ is far smaller for larger $L_1$.
\end{proof}
\section{The smoothed explicit formula}\label{sec:explicit}

We use the smoothing of B\"uthe \cite{Buthe2016}, with our own proof of the smoothing step (Lemma~\ref{lem:smooth}).
For $c > 0$ and $0 < \eps \le 10^{-4}$, put
\[
\eta_{c,\eps}(\tau) = \frac{c}{\eps\sinh c}\, I_0\!\Big(c\sqrt{1 - (\tau/\eps)^2}\Big)\mathbf 1_{|\tau|\le\eps},
\qquad
\ell_{c,\eps}(z) = \frac{c}{\sinh c}\,\frac{\sinh\sqrt{c^2 - (\eps z)^2}}{\sqrt{c^2-(\eps z)^2}} .
\]
Here $I_0$ is the modified Bessel function.
Then $\int\eta_{c,\eps} = 2$, and the Fourier transform satisfies $\int \eta_{c,\eps}(\tau)e^{-i z\tau}d\tau = 2\ell_{c,\eps}(z)$.
Logan's function $\ell_{c,\eps}$ is entire and even, and $\ell_{c,\eps}(0) = 1$.
Let $\lambda = \int\eta_{c,\eps}(\tau)e^{\tau/2}d\tau = 2\ell_{c,\eps}(i/2) \ge 2$, and let
\begin{equation}\label{eq:phi}
\phi(\tau) = \eta_{c,\eps}(\tau)e^{\tau/2}/\lambda, \qquad \psi_{c,\eps}(x) = \int \psi_0(xe^{-\tau})\,\phi(\tau)\,d\tau,
\end{equation}
where $\psi_0$ is the normalised Chebyshev function.
Thus $\phi$ is a probability density on $[-\eps,\eps]$, and $\psi_{c,\eps}(x)$ is a weighted average of $\psi_0$ over $[xe^{-\eps}, xe^{\eps}]$.
It coincides with B\"uthe's modified Chebyshev function \cite[(3.2)--(3.4)]{Buthe2018}.

\subsection{Smoothing}
\begin{lemma}\label{lem:smooth}
Let $x \ge 10^{18}$, let $c\in\{30,31,\dots,110\}$ be an integer, $0 < \eps \le 10^{-4}$, and $\lambda_s = \log(10\sqrt c/\eps) < \log x - 1$.
Put $\nu_c := \int_0^1 u\,\eta_{c,1}(u)\,du$. Then
\begin{multline*}
|\psi(x) - \psi_{c,\eps}(x)| \le \frac{\log x + \lambda_s}{\log x - \lambda_s}\cdot 2\eps\, e^{\eps/2}\Big(\frac{\nu_c}{2} + \frac{I_0(c)}{200\sinh c}\Big)x \\ + a(x),
\end{multline*}
where $a(x) = e^{2\eps}\log(e^{\eps}x)\big(2.01\eps\sqrt x + \tfrac12\log\log(2x^2)\big)$.
Moreover $2\sqrt c\,e^{\eps/2}\big(\nu_c/2 + I_0(c)/(200\sinh c)\big) \le 0.80319 < 0.7946\cdot 1.02$ for these $c$ (the scan uses only integer $c$).
\end{lemma}

\begin{proof}
Since $\phi$ is a probability density and $\psi$ is non-decreasing,
\[
\psi(x) - \psi_{c,\eps}(x) = \int (\psi(x) - \psi_0(xe^{-\tau}))\phi(\tau)\,d\tau \le \int_0^{\eps}\big(\psi(x) - \psi((xe^{-\tau})^-)\big)\phi(\tau)\,d\tau .
\]
For $\tau < 0$ the integrand is $\le 0$, because $\psi_0(y) \ge \psi(x)$ for $y > x$.

For $3 \le Y \le X/2$ the Brun--Titchmarsh inequality of Montgomery and Vaughan, $\pi(X)-\pi(X-Y) \le 2Y/\log Y$ \cite{MV73}, together with the trivial count of prime powers as in \cite[(3.9)]{Buthe2018}, gives
\[
\psi(X) - \psi(X - Y) \le \log X\Big(\frac{2Y}{\log Y} + 0.6\frac{Y}{\sqrt X} + \log\log X^2\Big).
\]
Put $Y = x(1-e^{-\tau}) \le x\tau$.
For $\tau \ge \tau_0 := \eps/(10\sqrt c)$ we have $\log Y \ge \log x - \lambda_s$.
For $\tau < \tau_0$ we bound the integrand by its value at $\tau_0$, using monotonicity.

The mass of $\phi$ below $\tau_0$ is at most $\tau_0\max\phi \le \tau_0 e^{\eps/2}c I_0(c)/(2\eps\sinh c)$.
Also $\int_0^\eps\tau\phi \le e^{\eps/2}\eps\nu_c/2$, since $\lambda \ge 2$.
Combining these gives the main term, with $\log x$ in the numerator.

The lower bound $\psi_{c,\eps}(x) - \psi(x) \le \int_{-\eps}^0 (\psi_0(xe^{|\tau|}) - \psi(x))\phi$ is treated in the same way.
Here the intervals end at $xe^{\eps}$, which produces the factor $\log(e^\eps x) \le \log x + \lambda_s$.
The prime-power and $\log\log$ terms give $a(x)$.

The final numerical inequality is certified in ball arithmetic by \code{certify/smoothing\_const.py}, using \code{acb.integral} for $\nu_c$.
The maximum over $c$ is attained at $c = 110$.
\end{proof}

\begin{remark}\label{rem:buthe}
\cite[Prop.~4]{Buthe2016} (quoted as \cite[Prop.~1]{Buthe2018}) states the main term as $2\eps x|\nu_c(\alpha)|\log(e^\eps x)/\log B$, with $\nu_c(t) = \int_{-\infty}^t\mu_c$ and $\mu_c(t) = -\int_{-\infty}^t\eta_{c,1}$ for $t<0$.
With $\eta_{c,1}$ normalised as above ($\int\eta_{c,1} = 2$), one has $|\nu_c(0)| = \nu_c \sim \sqrt{2/(\pi c)}$, whereas \cite[Remark~1]{Buthe2018} quotes $|\nu_c(0)|\sim(2\pi c)^{-1/2}$, which corresponds to a density of mass $1$.
We have not been able to reconcile the two normalisations.
Our own proof of the smoothing step (Lemma~\ref{lem:smooth}) does not depend on this.
Its main constant, $0.80/\sqrt c$, is half of what $2|\nu_c(0)|$ would give with the stated normalisation.
\end{remark}

\subsection{Explicit formula and truncation}
\begin{proposition}[{\cite[Prop.~2]{Buthe2018}, \cite[Props.~2,3]{Buthe2016}}]\label{prop:explicit}
Let $x \ge 10$, $0 < \eps \le 10^{-4}$ and $c \ge 3$. Then
\[
x - \psi_{c,\eps}(x) = \frac{1}{\ell_{c,\eps}(i/2)}\sum_\rho \ell_{c,\eps}\!\Big(\frac{\rho - 1/2}{i}\Big)\frac{x^\rho}{\rho} + \Theta(2),
\]
where $\Theta(2)$ denotes a quantity of absolute value at most $2$, and
\[
\sum_{|\Ip\rho| > c/\eps}\Big|\frac{\ell_{c,\eps}((\rho-\frac12)/i)}{\ell_{c,\eps}(i/2)}\frac{x^\rho}{\rho}\Big| \le 0.16\,\frac{x+1}{\sinh c}\,e^{0.71\sqrt{c\eps}}\log(3c)\log\frac c\eps .
\]
\end{proposition}

\begin{lemma}\label{lem:logancomplex}
Let $\rho = \beta + i\gamma$ with $0 < \beta < 1$ and $|\eps\gamma| \le c$. Then
\[
\Big|\ell_{c,\eps}\Big(\frac{\rho - 1/2}{i}\Big)\Big| \le e^{\eps(c+1)/6}\,\ell_{c,\eps}(\gamma),
\qquad \ell_{c,\eps}(i/2) \ge 1 .
\]
\end{lemma}

\begin{proof}
Write $\ell_{c,\eps}(z) = \frac{c}{\sinh c}\,G(c^2 - \eps^2z^2)$ with $G(w) = \sum_{n\ge0}w^n/(2n+1)!$.
For $z = \gamma - i(\beta-\frac12)$ we have
$|c^2 - \eps^2 z^2| \le (c^2 - \eps^2\gamma^2) + \eps^2/4 + \eps^2|\gamma| \le (c^2 - \eps^2\gamma^2) + \eps(c+1)$.
Since the coefficients of $G$ are positive, $|G(w)| \le G(|w|)$.
The bound $\frac{d}{dw}\log G(w) \le 1/6$ for $w \ge 0$ then gives the first claim.
The second holds because $G$ is increasing.
\end{proof}

\begin{lemma}\label{lem:sumrecip}
For $T \ge 10^4$, $\sum_{0<\gamma\le T} 1/\gamma \le \frac12\log^2 T$, the sum running over zeros with multiplicity.
\end{lemma}

\begin{proof}
Write the sum as a Stieltjes integral $\int_{14}^{T} t^{-1}dN(t) = N(T)/T + \int_{14}^T N(t)t^{-2}dt$.
Insert the explicit Riemann--von Mangoldt bound $|N(t) - \frac t{2\pi}\log\frac t{2\pi e} - \frac78| \le 0.1038\log t + 0.2573\log\log t + 9.3675$ for $t \ge e$ \cite{HSW}.
This gives $\sum_{0<\gamma\le T}1/\gamma \le \frac1{4\pi}\log^2T + \log T + 2$, which is below $\frac12\log^2T$ for $T \ge 10^4$.
At $T = H_0$ the bound is $96.5$, against $\frac12\log^2H_0 = 412.7$.
\end{proof}

\subsection{The bound for $\psi$ on an interval}
For $s \in [0.6, 1)$, $H = H_0 = 3\cdot 10^{12}$ and $T_1 = c/\eps$, put
\[
S(s) := \sum_{\substack{H < \gamma \le T_1\\ \beta > s}} \frac{\ell_{c,\eps}(\gamma)}{\gamma}\,e^{\eps(c+1)/6}.
\]

\begin{proposition}\label{prop:master}
Let $43 \le L_1 < L_2$, $c \in [30,110]$ and $\eps \le 10^{-4}$, with $\lambda_s < L_1 - 1$. Let $T_1 = c/\eps$.
Assume every zero $\beta+i\gamma$ with $H_0<\gamma \le T_1$ has $\beta \le s_{\max} := 1 - \lambda^*(\log T_1)/\log T_1$, where $\lambda^*$ is as in \S\ref{sec:zfr} (this holds by Theorem~\ref{thm:H}).
Then for every $x$ with $\log x\in[L_1,L_2]$,
\begin{multline*}
\frac{|\psi(x)-x|}{x} \le \frac{L_2+\lambda_s}{L_1-\lambda_s}\cdot\frac{0.7946\cdot 1.02\,\eps}{\sqrt c}
+ \frac{0.32\,e^{0.71\sqrt{c\eps}}\log(3c)\log T_1}{\sinh c} + e^{-L_1/2}\log^2 H\\
+ 2\Big(L_2\int_{0.6}^{s_{\max}} e^{(s-1)L_1}S(s)\,ds + e^{-0.4L_1}\big(S(0.6) + \log^2 T_1\big)\Big) + \frac{a(x)+2}{x}.
\end{multline*}
\end{proposition}

\begin{proof}
Combine Lemma~\ref{lem:smooth} with Proposition~\ref{prop:explicit}.
The zeros with $|\gamma| > T_1$ contribute the second term, since $0.16(x+1) \le 0.32x$.

Zeros with $|\gamma| \le H$ lie on the critical line \cite{PlattTrudgianRH}.
By Lemmas~\ref{lem:logancomplex} and~\ref{lem:sumrecip} they contribute at most $x^{1/2}\cdot2\cdot\frac12\log^2H$.

Zeros with $H < |\gamma| \le T_1$ come in pairs $\rho,\bar\rho$, which gives the factor $2$.
For those with $\beta > 0.6$, use the identity
\[
\sum_{\beta>0.6} x^{\beta-1}f(\gamma) = x^{-0.4}S(0.6) + \log x\int_{0.6}^1 x^{s-1}S(s)\,ds,
\]
together with $x^{s-1} \le e^{(s-1)L_1}$ and $\log x \le L_2$; also $S(s) = 0$ for $s \ge s_{\max}$.

All remaining zeros in this height range have $\beta \le 0.6$.
These include the reflections $1-\bar\rho$ of zeros off the line, and the zeros on the line.
They contribute at most $x^{-0.4}e^{\eps(c+1)/6}\sum_{H<\gamma\le T_1}1/\gamma \le x^{-0.4}\log^2T_1$; here Lemma~\ref{lem:sumrecip} gives $\frac12\log^2T_1$, which leaves ample room for the factor $e^{\eps(c+1)/6} \le 1 + 10^{-9}$.
\end{proof}

\begin{lemma}[Stieltjes bound]\label{lem:stieltjes}
Let $f(t) = \ell_{c,\eps}(t)/t$, which is decreasing on $(0, T_1]$.
Let $H = t_0 < t_1 < \dots < t_n = T_1$, and let $N^+(s,t)$ be any upper bound for $N(s,t) := \#\{\rho: \beta > s,\ 0<\gamma\le t\}$.
Suppose $N^+(s,t_{i+1})$ is replaced by $0$ whenever $s \ge 1 - \lambda^*(\log t_{i+1})/\log t_{i+1}$. Then
\[
S(s)e^{-\eps(c+1)/6} \le \sum_{i<n} N^+(s, t_{i+1})(f(t_i) - f(t_{i+1})) + N^+(s, T_1)f(T_1).
\]
\end{lemma}

\begin{proof}
This is summation by parts.
Zeros with $\gamma \le H$ do not satisfy $\beta > s \ge 0.6$, so they do not contribute.
The zero-free region is used at the top of each block; this is valid because $\lambda^*(l)/l$ is non-increasing in $l$.
\end{proof}

In the computation, $\int_{0.6}^{s_{\max}}e^{(s-1)L_1}S(s)\,ds$ is bounded by $\sum_j S(s_j)\int_{s_j}^{s_{j+1}}e^{(s-1)L_1}ds$, which is valid because $S$ is non-increasing.
The grid consists of $61$ points up to $s_{\max} - 0.02$, followed by steps $1/(40L_1)$.
The function $N^+$ is the minimum of the density bounds of \S\ref{sec:density}.
\section{A height-dependent zero-free region}\label{sec:zfr}

\subsection{The inequality of Bellotti, Trudgian and Yang}
We recall the set-up of \cite{BTY} and keep its notation.
\begin{itemize}
\item $\theta = 1.1338$, and $w$ is Heath-Brown's function \cite[Def.~1]{BTY}, supported on $[0,D]$ with $D = 2\theta\cot\theta$.
  Its Laplace transform is $W(z) = \int_0^D e^{-zu}w(u)\,du$, and $w_0 = w(0) = 5.672787598\ldots$.
\item $\kappa_m$ ($0 \le m \le 6$) are the coefficients in \cite[(18)]{BTY}, with $\kappa = \sum\kappa_m = 433/859$.
\item $P(x) = \sum_{k\le 16}a_k\cos kx \ge 0$ is the trigonometric polynomial of \cite[\S4]{BTY}, with $a_0 = 1$, $a_1 = 865534/497079$ and $a = \sum_{k\ge1}a_k = 2919857/828465$.
\item $K = 16$ and $T_0 = 10^{10}$.
\end{itemize}
For $\eta > 0$ and $\frac12 < \sigma < 1$, put
\[
f(u) = \eta w(\eta u)\sum_{m}\kappa_m e^{-(2\sigma-1)mu}, \qquad F(z) = \sum_m \kappa_m W\big((z + (2\sigma-1)m)/\eta\big).
\]
In particular $f(0) = \kappa\eta w_0$.

Let $\rho_0 = \beta_0 + it$ be a zero with $t \ge H_0$, and put $\eta = 1-\beta_0$ and $\lambda_0 = \eta\log t$.
Suppose $\zeta$ has no zeros with $\beta \ge 1 - A/\log\gamma$ for $H_0 \le \gamma \le Kt + T_0$, and put $\sigma = 1 - A/\log(Kt+T_0)$.
Assume $\eta \le \eta_0 \le 0.0145$ and $\sigma \ge \sigma_0 \ge 0.9930$.
The argument of \cite[Lemma~15]{BTY} then gives
\begin{multline}\label{eq:bty}
c_{10}\,\eta w_0 \le \eta w_0\Big(\frac{a\kappa}2\log t - G\Big) + a_0F(\sigma-1) - a_1\big(F(\sigma-1+\eta) + F(\sigma-\eta)\big) - Z\\ + 10^{-10}\eta + 2\cdot10^{-7} + E_2\lambda_0\eta^2 + E_3\eta^3 .
\end{multline}
Here $Z \ge 0$ is the part of the zero sum that is discarded.
The constants $c_{10}, G, E_2, E_3$ depend on $(\eta_0,\sigma_0)$ and are given in \S\ref{ss:constants}.
A zero with these parameters is impossible as soon as the following margin is positive:
\begin{multline}\label{eq:margin1}
\mathcal M_1 := a_1\big(F(\sigma-1+\eta)+F(\sigma-\eta)\big)\\ - \Big[\eta w_0\big(\tfrac{a\kappa}2\log t - G - c_{10}\big) + a_0F(\sigma-1) + 10^{-10}\eta + 2\cdot 10^{-7} + E_2\lambda_0\eta^2 + E_3\eta^3\Big] .
\end{multline}

\begin{remark}\label{rem:btyf0}
In \cite{BTY} the Lemma-10 term enters the final inequality as $-0.1186\,w(0)\eta$, i.e.\ with $f(0) = \eta w(0)$.
By their Definition~2, $f(0) = \kappa\eta w(0)$ with $\kappa = 0.504$.
With the correct value, the Lemma-10 gain $0.23545\,\kappa f(0)$ equals $0.0598\,\eta w_0$, not $0.1186\,\eta w_0$.
Evaluating their inequality exactly at the critical height $t = e^{76.47}$ with this correction still gives $\eta\log t \ge 0.2054 > 1/4.896$.
So their Theorem~1 is unaffected.
(We also note the misprint $F(2\sigma-1+\beta_0)$ in their (25), which should read $F(\sigma - 1 + \beta_0)$, as in (47).)
\end{remark}

\subsection{The constants for general $(\eta_0,\sigma_0)$}\label{ss:constants}
The constants of \cite{BTY} are computed for $\eta \le 0.0071093$ and $\sigma \ge 0.9935164$.
We need $\eta$ up to about $0.0125$ and $\sigma$ down to $0.9930$, so we re-derive each $(\eta_0,\sigma_0)$-dependent ingredient.

\begin{lemma}[$W_0$ bound; replaces {\cite[Lemma~4]{BTY}}]\label{lem:W0}
Put $W_0(z) = W(z) - w_0/z$. For $\Rp z \ge \nu$ and $|z| \ge r > 0$,
\[
|W_0(z)| \le C_w(\nu,r)|z|^{-3}, \qquad C_w(\nu,r) = |w''(0)| + \frac{|w'''(0)| + e^{-\nu D}|w'''(D)| + \int_0^D|w''''(u)|e^{-\nu u}du}{r} .
\]
Here $|w''(0)| = 50.06970149\ldots$ and $|w'''(0)| = |w'''(D)| = 117.26477652\ldots$.
\end{lemma}

\begin{proof}
We have $w(D) = w'(D) = w''(D) = 0$ and $w'(0) = 0$.
Integrating by parts four times gives
$W(z) = \frac{w_0}{z} + \frac{w''(0)}{z^3} + \frac{w'''(0) - e^{-zD}w'''(D)}{z^4} + z^{-4}\int_0^De^{-zu}w''''(u)\,du$.
The derivative values are exact enclosures from Taylor series in \code{arb}.
The integral is certified on an integer grid of $\nu \in [30,200]$, using that it is decreasing in $\nu$.
Sign changes of $w''''$ are isolated, and each piece is integrated with \code{acb.integral} (\code{rep/consts\_cert.py}).
\end{proof}

For comparison, $C_w(69.4, 69.4) = 51.80$, while \cite{BTY} use $C < 52$ at this point.

\begin{lemma}[non-negativity; replaces {\cite[Lemma~5, (21)]{BTY}}]\label{lem:L5}
Let $\sigma\in[0.993,1)$ and $\eta \le \eta_0 \le 0.0145$. Then $\Rp F(z) + \Rp F(2\sigma - 1 - \bar z) \ge 0$ for $0 \le \Rp z \le 2\sigma-1$.
\end{lemma}

\begin{proof}
As in \cite{BTY}, it suffices to treat $\Rp z = 0$.
For fixed $\sigma$ every $\delta_m$ equals $(2\sigma-1)m$ exactly.
Rescaling $y \mapsto y/(2\sigma-1)$ reduces the condition to
\[
B(y) := \sum_{m=1}^{7} b_m\frac{m - \sgn(b_m)\eps_0/m}{m^2+y^2} > 0 \quad (y\in\mathbb R), \qquad \eps_0 = \frac{C_w\eta_0^2}{w_0(2\sigma-1)^2}.
\]
Here $(1+x)\sum\kappa_mx^m = \sum b_mx^m$.
This reduction is independent of $\sigma$.
For $\eta_0 \le 0.0145$ and $\sigma \ge 0.993$ we have $\eps_0 \le 1/500$ (with $C_w \le 51.9$).
For $\eps_0 = 1/500$, $B$ is a rational function whose numerator has no real roots and whose value at $0$ is positive; this is checked by exact root isolation (\code{rep/lemma5\_cert.py}).
\end{proof}

\begin{remark}\label{rem:L5}
\cite{BTY} prove (21) uniformly in $\sigma$, replacing $\delta_m$ by $c_0m$ with $c_0 = 151/153 \le 2\sigma_0-1$.
This requires $\sigma \ge 0.993464$.
For $c_0 = 0.9864$, which corresponds to the smaller values of $\sigma$ needed below, their test function $B$ takes negative values, already at $\eps_0 = 1/2000$.
The per-$\sigma$ version above has no such restriction.
\end{remark}

\begin{lemma}[error terms; replaces {\cite[Lemmas~8, 9 and (48)]{BTY}}]\label{lem:L89}
Let $x_0 = \sigma_0 - \frac12 \ge 0.493$, $C_1 = C_w(x_0/\eta_0, x_0/\eta_0)$ and $C_2 = C_w(\sigma_0/\eta_0,\sigma_0/\eta_0)$.
Let $\Sigma_\kappa = \sum|\kappa_m|$. Then $\sum_k a_k|E_k| \le E_2\lambda_0\eta^2 + E_3\eta^3$, where
\[
E_2 = a\,\Sigma_\kappa\frac{C_1}{\pi x_0^2}, \qquad
E_3 = a_0\Sigma_\kappa\Big(\frac{C_1}{2\pi}\cdot 23.6497 + \frac{C_2}{x_0^3}\Big) + a\,\Sigma_\kappa\frac{C_2}{x_0^3} .
\]
\end{lemma}

\begin{proof}
Follow \cite[Lemma~8]{BTY}, with two changes.
First, in their (33) the inequality $\eta\log t \le A_0$ is replaced by $\eta\log t = \lambda_0$.
Second, the integral (32) is evaluated with the true digamma function:
$\int_{\mathbb R}|\Rp\psi(\frac14 + \frac{iy}2)|(x_0^2+y^2)^{-3/2}dy \le 23.6497$ at $x_0 = 0.4930$, certified (\code{rep/consts\_cert.py}).
This integral is decreasing in $x_0$.
\end{proof}

\begin{remark}
The majorant $U(y)$ of \cite[(31)]{BTY} integrates to $28.3$ against $(x_0^2+y^2)^{-3/2}$, not to less than $24$ as stated in (32).
The true digamma gives $23.59$, so (32) is correct, but its proof as written does not establish it.
At $(\eta_0,\sigma_0) = (0.0071093, 0.9935164)$, Lemma~\ref{lem:L89} gives $185.4\eta^2 + 7972\eta^3$, against $187\eta^2 + 8013\eta^3$ in \cite[(48)]{BTY}.
\end{remark}

\begin{lemma}[{\cite[Lemmas~3, 10, 11]{BTY}}, re-derived]\label{lem:L1011}
\begin{enumerate}
\item $f \ge 0$ on $[0,\infty)$.
\item $f(u) \ge \kappa f(0)$ for $0 \le u \le u_0(\eta_0,\sigma_0)$; for example $u_0 = 59.15$, $49.46$ and $33.61$ at $\eta_0 = 0.0071$, $0.0085$ and $0.0125$.
\item $\sum_n\Lambda(n)n^{-\sigma}f(\log n)P(t\log n) \ge c_{10}\,\eta w_0$, where $c_{10} = \kappa^2\sum_{p:\ 15\log p\le u_0}m_p\log p$.
\end{enumerate}
Here $m_p$ are the certified lower bounds for $\min_x G_p(1,x)$ in \cite[Table~1]{BTY}, except that $m_5 = 0.0073564$ (Remark~\ref{rem:m5}).
\end{lemma}

\begin{proof}
(1) $w \ge 0$ on $[0,D]$; this is certified by ball bisection on $[0,D-0.02]$ and by a Taylor bound near $D$, where $w'''(D) < 0$.
The polynomial $\sum\kappa_mx^m$ has no root in $[0,1]$, by exact root isolation (\code{rep/f\_nonneg\_cert.py}).

(2) We have $g(u) := w(\eta_0u)\sum\kappa_me^{-d_mu} - \kappa w_0 \ge 0$, with $d_m$ as in \cite[Lemma~11]{BTY}.
This is certified by ball bisection, using $g(0) = 0$ and $g'(0) = -w_0\sum\kappa_md_m \ge 1.19$ near $u = 0$ (\code{rep/lemma11\_cert.py}).

(3) As in \cite[Lemma~10]{BTY}, the terms with $n > e^{u_0}$ are dropped, which is valid by (1) and $P \ge 0$.
So are all primes with $15\log p > u_0$, and the powers $p^m$ with $m > 15$.
The bound $G_p(\sigma,x) \ge G_p(1,x)$ needs no computation: $G_p(\sigma,x) = \sum_{m\le15}p^{-m\sigma}P(mx)$, so $\partial_\sigma G_p = -\log p\sum_m m\,p^{-m\sigma}P(mx) \le 0$.
The values $m_p$ are re-certified by ball bisection with a mean-value enclosure (\code{rep/mp\_cert.py}).
\end{proof}

\begin{remark}\label{rem:m5}
For $p = 5$, \cite[Table~1]{BTY} lists $m_5 = 0.00735709$.
However $G_5(1,x)$ attains $0.0073564919\ldots$ at $x = 1.6495242\ldots$, as confirmed independently in $30$-digit arithmetic.
All other entries of the table are confirmed.
The effect on their constant $0.23545$ is $-10^{-6}$.
\end{remark}

\begin{lemma}[{\cite[Lemma~12]{BTY}}]\label{lem:L12}
$\sum_k a_k\sum_m\kappa_m\big(-\frac12\log\pi + \frac12\Rp\frac{\Gamma'}{\Gamma}(\frac{s_k+\delta_m}2+1)\big) \le \frac{a\kappa}2\log t - G(\sigma_0)$, where $G(\sigma_0)$ is computed by their formula in ball arithmetic.
We have $G(0.9930) = 1.56839$ and $G(0.9935164) = 1.56888$; \cite{BTY} give $1.568$.
\end{lemma}

Lemma~13 of \cite{BTY} ($10^{-10}\eta$) and the tail estimates in their Lemma~6 hold unchanged for $\eta \le 0.015$ and $\log t \le 45$.
Each tail term in their Lemma~6 is at most $3\cdot 10^{-9}$, so the total stays below $10^{-7}$; we use $2\cdot10^{-7}$.

\subsection{Exact evaluation of $W$}
Since $w$ is a combination of $u\cos(ku)$, $\cos(ku)$, $\sin(\phi-ku)$, $u$ and $1$ on $[0,D]$ (with $k = \tan\theta$), $W(z)$ has a closed form.
In terms of $J_0(s) = (1-e^{-sD})/s$ and $J_1(s) = (1 - e^{-sD}(1+sD))/s^2$ it reads
\[
W(z) = \sec^2\theta\Big[\sec^2\theta\cdot\tfrac12\sum_\pm\big(\theta\cot\theta\, J_0 - \tfrac12J_1\big)(z\mp ik) + (2\theta\cot\theta-2)J_0(z) - J_1(z) + \frac{\Sigma_{2\theta}(z)}{\sin2\theta} - \frac{2\Sigma_\theta(z)}{\sin\theta}\Big],
\]
where $\Sigma_\phi(z) = (e^{i\phi}J_0(z+ik) - e^{-i\phi}J_0(z-ik))/(2i)$.
It agrees with quadrature to $10^{-15}$ (\code{rep/Wexact.py}).
In certification we evaluate it in \code{acb}, with the second-order enclosure
\[
W(z) \in W(z_c) + W'(z_c)(z - z_c) + B\big(0, \tfrac12 r^2\sup|W''|\big), \qquad r = |z-z_c| .
\]
Here $W'$ has the same closed form with $J_0\to-J_1$ and $J_1\to-J_2$.
Integrating $u^2w$ by parts three times gives $\sup|W''| \le \min\big(\int u^2we^{-\nu u},\ (2w_0 + \int|(u^2w)'''|e^{-\nu u})/|z|^3\big)$.
Both integrals are bounded by sums of ball suprema over $256$ subintervals.

\subsection{The zero-free region}
\begin{theorem}\label{thm:H}
Let $\lambda^*(l)$ be the function given by the table below on the grid $l \in \{28, 28.25, \dots, 43\}$, extended by $\lambda^*(l) := \lambda^*(\lceil l\rceil_{0.25})$ for $l \le 43$ and by $\lambda^*(l) := 1/4.896$ for $l > 43$.
Then every zero $\beta+i\gamma$ of $\zeta$ with $\gamma > H_0$ satisfies $(1-\beta)\log\gamma \ge \lambda^*(\log\gamma)$; for $3\le\gamma\le H_0$ all zeros have $\beta = 1/2$ \cite{PlattTrudgianRH}.
\begin{center}\small
\begin{tabular}{c|cccccc}
\toprule
$l$ & 28 & 29 & 30 & 31 & 32 & 33\\
$\lambda^*(l)$ & 0.21424 & 0.21384 & 0.21345 & 0.21308 & 0.21272 & 0.21239\\
\midrule
$l$ & 34 & 35 & 36 & 37 & 38 & 40\\
$\lambda^*(l)$ & 0.21206 & 0.21175 & 0.21145 & 0.21117 & 0.21084 & 0.21033\\
\bottomrule
\end{tabular}
\end{center}
The full table is \code{rep/lamstar2c.json}.
\end{theorem}

\begin{proof}
For $\gamma \le H_0$ the statement follows from \cite{PlattTrudgianRH}.
For $l > 43$ it is \cite[Thm.~1]{BTY}.
For the remaining range we run a four-step chain, starting from $\lambda_0 \equiv 1/4.896$.

At step $j$ the hypothesis is the conclusion of step $j-1$.
Since $\lambda_{j-1}(l)/l$ is non-increasing, this hypothesis gives the zero-free input of Lemma~6 of \cite{BTY} with $A = \lambda_{j-1}(\log(Kt+T_0))$.
In every height cell $(L - 0.25, L]$ we certify $\mathcal M_1 > 0$ in ball arithmetic.
This is done for all $\log t$ in the cell and all $\lambda_0 = \eta\log t \in [\lambda_{j-1}(L), \lambda_j(L)]$, where $\lambda_j(L)$ is the float threshold minus $2\cdot10^{-5}$, by adaptive bisection of the box.
On each box we take $\sigma$ as a ball, with $A$ at the top of the box, which is valid because a smaller $A$ is a weaker hypothesis.
The constants are taken at $\eta_0 = \lceil \lambda_j/l_{\min}\rceil_{10^{-4}} + 10^{-4}$ and at $\sigma_0$ rounded down.
Then $\lambda_j(L) \ge \lambda_{j-1}(L)$, and a zero with $\lambda_{j-1}(L) \le (1-\beta)\log\gamma < \lambda_j(L)$ would give $\mathcal M_1 > 0$, contradicting \eqref{eq:bty}.

We use the table produced by the fourth step; it differs from the third by at most $10^{-5}$ in every entry.
Program: \code{rep/lamcert.py}.
\end{proof}

At $\log t = 29$ the theorem amounts to the zero-free region $\sigma > 1 - 1/(4.676\log t)$.
\section{Repulsion of zeros near the $1$-line}\label{sec:rep}

In \cite[Lemma~6]{BTY} the zero sum $\Rp\sum_\rho F(s_1-\rho)$ is bounded below by the pair $\rho_0, 1-\bar\rho_0$.
Every other zero $\rho$ is paired with its reflection $1-\bar\rho$, and each such pair contributes $\ge 0$ by (21) (our Lemma~\ref{lem:L5}).
Instead of discarding one more pair, we keep it.

\begin{lemma}\label{lem:twozero}
In the situation of \eqref{eq:bty}, suppose $\rho_1 = \beta_1 + i(t+d)$ is a second zero with $\beta_1 > 1/2$.
Put $\eta_1 = 1-\beta_1$, and assume $\lambda^*$-admissibility, i.e.\ $\sigma - \beta_1 > 0$. Then \eqref{eq:bty} holds with
\[
Z \ge a_1\Rp\big(F(\sigma - 1 + \eta_1 - id) + F(\sigma - \eta_1 - id)\big) .
\]
\end{lemma}

\begin{proof}
The pair $\{\rho_1, 1-\bar\rho_1\}$ contributes exactly $\Rp F(s_1 - \rho_1) + \Rp F(s_1 - (1-\bar\rho_1))$ to the zero sum, and this equals the displayed quantity.
All other pairs are $\ge 0$ by Lemma~\ref{lem:L5}, since $0 \le \Rp(s_1-\rho) \le 2\sigma - 1$ for every zero below height $Kt+T_0$.
The term $s_1$ enters with the coefficient $a_1$.
The case $d < 0$ is identical, because $F(\bar z) = \overline{F(z)}$, so only $|d|$ matters.
\end{proof}

A pair of zeros at heights $t, t+d$ is therefore impossible as soon as
\[
\mathcal M_2 := \mathcal M_1(\eta_0) + a_1\Rp\big(F(\sigma-1+\eta_1-id) + F(\sigma-\eta_1-id)\big) > 0
\]
for at least one of the two choices of primary zero.
Here $F$ is formed with the primary $\eta$.
The reflected term behaves like $\Rp F(\sigma - \eta_1 - id) \approx \eta w_0(\sigma-\eta_1)/((\sigma-\eta_1)^2 + d^2)$, which is of size about $0.04$ for $d \lesssim 0.05$ and $\eta, \eta_1 \approx 0.0075$.
This is what produces the repulsion.

\begin{theorem}\label{thm:rep}
For integers $29 \le l_k \le 37$ and $k$ with $\lambda^*(l_k) < k\cdot10^{-4}\cdot l_k \le 0.0125(l_k-1)$, define $d^*(l_k,k)$ by the table \code{rep/dstar2c\_all.json}.
For each $l_k$ the value is largest for the smallest admissible $k$ (for example $0.142$ at $l_k = 29$), and for the large values of $k$ it equals $0.058$ ($l_k = 29$), $0.056$ ($l_k = 30$) and $0.045$ ($l_k = 37$).
Suppose two zeros $\beta_j + i\gamma_j$ ($j = 0,1$) satisfy $H_0 < \gamma_j \le e^{l_k}$ and $1-\beta_j \le k\cdot 10^{-4}$.
Then $|\gamma_0 - \gamma_1| \ge d^*(l_k, k)$.
\end{theorem}

\begin{proof}
For a cell $(l_k,k)$ put $e_{\min} = \lambda^*(l_k)/l_k$ and $e_{\max} = k\cdot10^{-4}$; the variables $\eta_0,\eta_1$ below denote $1-\beta$ of the two zeros.
By Theorem~\ref{thm:H}, both zeros have $1-\beta_j \ge e_{\min}$, since $\lambda^*(l)/l$ is non-increasing.

Fix a sub-cell $(l-1, l]$ of heights, $l \le l_k$.
We certify $\mathcal M_2 > 0$ for all $(\log t, \eta_0, \eta_1, |d|) \in [l-1, l + 10^{-12}] \times [e_{\min},e_{\max}]^2 \times [0, d^*]$, in ball arithmetic by adaptive bisection of the box.
The width $10^{-12}$ covers the height $t+d$ of the second zero.
On each box we try both orderings of the pair.
As before, $\sigma$ is taken as a ball, with $A$ at the top of the box, and the constants at $(\eta_0,\sigma_0)$ rounded outward.
In this way all $k$ on a coarse grid are certified at $97\%$ of the float value of $d^*$ (\code{rep/rep\_cert.py}, \code{rep/repcert\_run.py}).
Intermediate $k$ inherit the value of the next certified $k$ above, which is valid because enlarging $e_{\max}$ only enlarges the box.

Finally, a zero counted at height $\le e^{l_k}$ may lie in any lower height cell.
So the table records $\min_{l\le l_k} d^*(l,k)$, over the cells $l$ in which $1-\beta \le k\cdot10^{-4}$ is compatible with Theorem~\ref{thm:H} (\code{rep/min\_fix.py}).
This minimum changes only the six cells $(l_k, 121)$ with $l_k \ge 32$, which become $0$.
\end{proof}

The first version of these certificates used a second-order enclosure for $W''$ that clamped $\nu$ from below at $-2$.
An independent audit found this.
After the fix, every one of the $131$ certified cells, and the chain of Theorem~\ref{thm:H}, was certified again with identical results.
\section{Zero-density estimates}\label{sec:density}

Let $N(\sigma,T) = \#\{\rho = \beta+i\gamma: \beta > \sigma,\ 0 < \gamma \le T\}$, counted with multiplicity.
In Proposition~\ref{prop:master}, $N(s,t)$ is replaced by $N^+(s,t)$, the minimum of the following bounds, each valid in its stated range.
\begin{enumerate}
\item[(N1)] $N(T) \le \frac T{2\pi}\log\frac T{2\pi e} + \log T + 1$ for $T \ge H_0$; this follows from \cite{HSW}.
\item[(N2)] For $\sigma \ge \sigma_0 \in\{0.60, 0.65, \dots, 0.99\}$ and $T \ge 3.061\cdot 10^{10}$: Kadiri--Lumley--Ng \cite[Table~1]{KLN}, in the form $N(\sigma,T) \le A(\log T)^{5-2\sigma_0}T^{\frac83(1-\sigma_0)} + B\log^2T$.
  This is valid because $\sigma \ge \sigma_0$, and the rows with $k < 1$ are used as $k = 1$.
\item[(N3)] For $\sigma \ge 0.9927$ and $T > H_0$: Bellotti's bounds $N(\sigma,T) \le CT^{B(1-\sigma)}$ \cite[Thm.~1.1, Table~2]{Bellotti}, uniformly and for the ranges $T\in(e^{l-1},e^l]$ ($29\le l\le41$).
\item[(N4)] For $1-\sigma \le 0.03$ and $28 < \log T \le 40$: Theorem~\ref{thm:density} below.
  It is this bound that matters in the binding range.
\end{enumerate}

\subsection{The log-free estimate}
We follow Bellotti \cite[\S\S3--4]{Bellotti}. In this section $L = \log T$ (not $\log x$), and $U = T^u$, $V = T^v$ ($v = \tau u$), $W = T^w$ and $X = T^x$.
\begin{itemize}
\item The Barban--Vehov weights are $\psi_d$, supported on $d\le V$, with $\Psi(n) = \sum_{d\mid n}\psi_d$, so that $\Psi(1) = 1$ and $\Psi(n) = 0$ for $1 < n \le U$.
\item The Selberg-type weights are $\theta_d$, supported on $d\le W$, with $\Theta(n) = \sum_{d\mid n}\theta_d$.
\end{itemize}
Let $\alpha = 1-\delta$.
\begin{itemize}
\item The plane is cut into rectangles of height $h = \kappa_0\delta$.
\item In each rectangle we choose one representative zero with $\beta \ge \alpha$, if there is one; let $J$ be the number of representatives.
\item Let $\tilde n$ be an upper bound for the number of zeros with $\beta\ge\alpha$ in one rectangle.
\end{itemize}

\begin{theorem}\label{thm:density}
Let $\delta \in ((k-1)10^{-4}, k\,10^{-4}]$ with $k \le 300$, and $\log T \le l$ with $l \in \{29, 29.5, 30, 30.5, 31, 32, \dots, 40\}$.
Let $(u,\tau,w,x)$ be admissible for the cell $(k,l)$ in the sense of Lemma~\ref{lem:D}.
Then
\[
N(1-\delta, T) \le N(1-\delta_k, e^l) \le \tilde n\,\frac{d_4d_5}{w}\cdot\frac{\mathcal C}{c_1^2 - c_5}\cdot\frac{e^{2x\delta_kl}}{\delta_kl},
\]
where $\delta_k = k\,10^{-4}$ and the constants $\tilde n, d_4, d_5, \mathcal C, c_1, c_5$ are those of Lemma~\ref{lem:D}, evaluated at $\delta_k$ and $L = l$.
(The first inequality is monotonicity of $N$; no rescaling within a cell is used.)
\end{theorem}

The parameters are chosen separately for each cell by minimising the bound numerically.
Any admissible choice gives a valid bound, so the optimisation needs no certification.
The bound is applied at the corner $(1-\delta_k, e^l)$ of the cell, which dominates the whole cell by monotonicity.
The resulting tables are \code{code/lemmaD\_tableGH0.01\_*.pkl}.

\begin{lemma}\label{lem:D}
Write $\alpha_k = 1 - \delta_k$ and $\nu = \lambda^*(l)/l$.
Call $(u,\tau,w,x)$ admissible if
\[
0.02 < w < u,\quad \tau > 1.05,\quad xL > uL - 2\log L,\quad uL \ge 40,\quad xL > \tau uL,
\]
\[
\delta_k > \nu,\quad c_1 > 0,\quad c_1^2 - c_5 \ge 0.2 .
\]
For admissible parameters the bound of Theorem~\ref{thm:density} holds with the following constants.
\begin{enumerate}
\item[(i)] $c_1 = e^{-L^2/U} - B_1 - 3.8\,e^{-L(L-2u)}\dfrac{L^2+1}{w(v-u)L^6}$, where $I_\Gamma = 5.922$ and
\[
B_1 = \frac{10.45\cdot 0.618\cdot L\cdot I_\Gamma}{2\pi}\,T^{(v+w)/2 + 1/6 + x(1/2-\alpha_k)} .
\]
\item[(ii)] $d_5 = 1.0061\big(1 + (1.333 + 4e^{-wL/2})/(wL)\big)$ \cite[Lemma~4.3]{Bellotti}.
\item[(iii)] $d_4$ is any upper bound for $\sum_n \Psi(n)^2n^{-1}(e^{-n/X} - e^{-nL^2/U})$; we take the bound of Theorem~\ref{thm:G}.
\item[(iv)] $c_5 = \dfrac{d_4d_5}{2\pi}\,\dfrac{e^{wL}}{wL}\Big(e^{Lx(3/2-2\alpha_k)} + e^{(Lu-2\log L)(3/2-2\alpha_k)}\Big)\,\mathcal J$, where $\mathcal J = 4.72\cdot 0.618\, e^{L/6}L$.
\item[(v)] $\mathcal C = \min_{\kappa_0}\Big[1.0067\big(1 + 2\delta_k\,\Lambda_{\mathrm{pair}}(\kappa_0\delta_k)\big)\Big] + 2\sqrt{2\pi}\cdot 2.4\cdot I'(\delta_k,l)$, with $\Lambda_{\mathrm{pair}}$ from Lemma~\ref{lem:pair} and $I'$ from Lemma~\ref{lem:crowd}.
\item[(vi)] $\tilde n = \max\big(1, \min_\eta\lfloor M(\eta,l)/\mathrm{cost}(h/2)\rfloor\big)$ with $h = \kappa_0\delta_k$, in the notation of Lemma~\ref{lem:pair}.
\end{enumerate}
\end{lemma}

\begin{proof}
We indicate only where the argument departs from \cite{Bellotti}.
The overall structure is \cite[(4.22)--(4.23)]{Bellotti}: $J^2c_1^2 \le J(c_2+c_3+c_4)X^{2-2\alpha}/(\delta\log T) + J^2c_5$, followed by $N(\alpha,T) \le \tilde nJ$.
Each change below replaces one of her constants by a smaller valid one.

\emph{(i) The term $B_1$.}
This replaces Bellotti's $44.66\,d_2(v+w)^4T^{(v+w)/2+x(1/2-\alpha_0)+27/164}\log^4T$ in \cite[(3.17)]{Bellotti}.
Write $a = ga'$ and $b = gb'$ with $(a',b') = 1$, and use $|\psi_a|, |\theta_b| \le 1$ and $\sum_{n\le N}n^{-1/2}\le 2\sqrt N$.
This gives
$|M(\tfrac12+it)| \le \sum_g g^{-1/2}\cdot 2(V/g)^{1/2}\cdot 2(W/g)^{1/2} < 4\zeta(3/2)\sqrt{VW} = 10.45\sqrt{VW}$,
a bound without logarithms.
We also use $|\zeta(\frac12+iu)| \le 0.618|u|^{1/6}\log|u|$ for $3 \le |u| \le e^{105}$ \cite{HPY}, and $\le 1.461$ for $|u| < 3$ \cite{Hiary}, together with
$(\gamma+|t|)^{1/6}\log(\gamma+|t|) \le T^{1/6}\log T\,(1+|t|)$.
Finally $I_\Gamma = \sup_{\sigma\in[-\frac12,\frac12-\alpha]}\int|\Gamma(\sigma+it)|(1+|t|)\,dt \le 5.922$ for $\delta \le 0.03$.
The last term of $c_1$ is Bellotti's bound for the tail $\sum_{n>U}\Psi\Theta n^{-\rho}e^{-nL^2/U}$ \cite[(3.19)ff.]{Bellotti}.

\emph{(iv) The term $\mathcal J$.}
This replaces $J_1+J_2+J_3+J_4 \le 0.241 + 5.921 + 286.077\,T^{27/164} + 0.241$ \cite[(4.17)--(4.19)]{Bellotti}.
We have $c := \frac32 - \beta_r - \beta_s \in [-\frac12, -0.44]$ for $\delta\le0.03$, and $|\zeta(\frac12 + iu)| \le \max(1.461, 0.618|u|^{1/6}\log|u|)$ for all $|u|\ge3$.
For $|u| > e^{105}$ this holds because Patel and Yang's $66.7|u|^{27/164}$ \cite{PY} is smaller there.
Moreover $|u| \le T + |t|$.
Hence the whole integral \cite[(4.16)]{Bellotti} is at most
\[
0.618\,T^{1/6}\log T\int_{\mathbb R}|\Gamma(c+it)|\Big(1+\frac{|t|}3\Big)^{1/6}\Big(1 + \frac{\log(1+|t|/3)}{\log3}\Big)dt \le 4.72\cdot 0.618\,T^{1/6}\log T .
\]
The integral equals $4.7136$ at $c = -0.44$ and increases with $c$.

\emph{(iii) The diagonal.} $d_4$ is the Barban--Vehov mean square of \cite[Lemma~4.2]{Bellotti}, which we bound by Theorem~\ref{thm:G}.
Bellotti's own bound \cite[Lemma~3.3, (3.6)]{Bellotti} gives $d_4 \approx 85$--$105$ at her parameters (her Table~5); Theorem~\ref{thm:G} gives about $1.5$ at ours.
Because Theorem~\ref{thm:G} bounds the whole sum, Bellotti's condition $u+v<x$ is not needed.

\emph{(v) The off-diagonal terms $c_2+c_3$.}
Bellotti assumes \cite[(4.11)]{Bellotti} that representatives of distinct rectangles are at least $\delta$ apart.
This need not hold: adjacent rectangles can have representatives arbitrarily close.
Lemma~\ref{lem:pair} makes no such assumption, and it also uses the repulsion of Theorem~\ref{thm:rep}.
The far term $c_4$ is replaced using Lemma~\ref{lem:crowd}.

\emph{(vi) The local count.} $\tilde n$ replaces $b_1\delta\log T + b_2$ \cite[Lemma~2.4]{Bellotti}.
\end{proof}

\subsection{The pair sum}
For $s_0 = 1 + \eta + i\gamma_r$, the Hadamard product and Stirling's formula give
\[
\sum_\rho\Rp\frac{1}{s_0-\rho} \le M(\eta, L) := \min\Big\{\tfrac12\log\tfrac{T}{2\pi} + \mathcal Z(\eta),\ \min_{c\in[0.5,2]}\Big(\tfrac12\log\tfrac T{2\pi} + \tfrac{c^2+1/2}{2c}\mathcal Z(\eta) + \tfrac{\log(T/\pi)}{8c}\Big)\Big\},
\]
where $\mathcal Z(\eta) = -\zeta'(1+\eta)/\zeta(1+\eta)$.
The second expression comes from positivity of $(c + \cos\phi)^2$.
Every term in the sum is positive.
A zero with $\beta \ge 1-\delta$ at vertical distance $d$ from $s_0$ contributes at least
$\mathrm{cost}(d) := \min_{a\in\{\eta, \eta+\delta\}} a/(a^2+d^2)$,
which is valid because $a\mapsto a/(a^2+d^2)$ is unimodal.

\begin{lemma}\label{lem:pair}
Fix a representative $r$ and $h = \kappa_0\delta$.
On each side of $r$, the representative of the $j$-th rectangle away lies at distance $d \in ((j-1)h, (j+1)h)$.
Its contribution to the Hal\'asz--Montgomery pair sum, relative to the diagonal, is at most
$b_j = g\big(\max(2\delta, (j-1)h, d^*)\big)/\max\big(2\delta, (j-1)h, d^*\big)$, where
$g(d) = \min\big(1, |\Gamma(1+2\delta+id)|/\Gamma(1+2\delta)\big)$ and $d^* = d^*(\lceil l\rceil, k)$ is the repulsion distance of Theorem~\ref{thm:rep}.
Its cost is at least $\mathrm{cost}((j+1)h)$.
Rectangles with $(j+1)h \le d^*$ contain no representative.
Then
\[
\sum_{s\ne r,\ |\gamma_s - \gamma_r| \le 1}\frac{|\Gamma(2-\rho_r-\bar\rho_s)|X^{2-\beta_r-\beta_s}}{1.0067\,X^{2\delta}} \le 2\delta\,\Lambda_{\mathrm{pair}}(h),
\]
where $\Lambda_{\mathrm{pair}}(h) := \min_\eta\max\big\{\sum_j b_jx_j : 0\le x_j\le 1 \ (\text{two per } j),\ \sum_j\mathrm{cost}((j+1)h)x_j \le M(\eta,L) - \mathrm{cost}(0)\big\}$.
\end{lemma}

\begin{proof}
We have $|\Gamma(z)| = |\Gamma(1+z)|/|z|$.
For $0\le\Rp z = s\le2\delta$, $|\Gamma(1+s+id)| \le \Gamma(1+s)|\Gamma(1+2\delta+id)|/\Gamma(1+2\delta)$, since the product $\prod_n(1+d^2/(1+s+n)^2)^{-1}$ is increasing in $s$.
Also $X^y/y$ is increasing on $[2\nu, 2\delta]$, which gives the factor $1/\max(2\delta, |\Delta|)$.
The crowding budget at $s_0 = 1+\eta+i\gamma_r$ bounds the total cost of all representatives near $\gamma_r$.
The maximum of the linear functional under this budget is the fractional knapsack value, computed greedily.
The repulsion theorem applies because all representatives counted in the cell $(k,l)$ have $1-\beta\le\delta_k$ and height $\le e^l$; for half-integer $l$ we use the integer cell above.
\end{proof}

\begin{lemma}[crowding]\label{lem:crowd}
For a window of height $\omega$ centred at $t \le T$ and any $\eta > 0$, the number of representatives in the window is at most
$n(\omega) = \min\big(\lfloor M'/\min(f_\omega(\eta), f_\omega(\eta+\delta))\rfloor, \lceil \omega/\delta\rceil + 1\big)$,
where $f_\omega(a) = a/(a^2 + \omega^2/4)$ and $M' = \frac12\log\frac T{2\pi} + \frac1\eta - 0.36$.
Consequently Bellotti's $c_4$ \cite[(4.13)--(4.14)]{Bellotti} may be computed with
$I' = \delta\sum_i n(\omega)\max_{[y_i, y_i+\omega]} g_\alpha$ in place of $I = \int_1^T g$, where $g_\alpha(y) = e^{y/6}y^{1.5-2\alpha}e^{-\pi y/2}$, $\omega = 1/4$, and the windows tile $[1,60]$.
\end{lemma}

\begin{proof}
This is the crowding argument above applied to $s = 1+\eta+it$.
For $0 < \eta \le 1.5$ we have $-\zeta'(1+\eta)/\zeta(1+\eta) \le 1/\eta - 0.36$; the minimum of $1/\eta + \zeta'/\zeta(1+\eta)$ on $(0, 1.5]$ is $0.378$.
Beyond $y = 60$ the function $g_\alpha$ is below $10^{-30}$.
The worst case is taken over each $\delta$-bin, and $\eta$ is optimised on a grid.
\end{proof}

\subsection{The Barban--Vehov mean square}
Let $\Delta = \log(V/U)$, $Y_1 = \log(X/V)$ and $Y_2 = \log(X/U)$.

\begin{theorem}\label{thm:G}
Let $\log U \ge 40$, $X > V$, $\tau \in [1.3, 2.5]$ and $\log X/\log V \in [1.1, 2.6]$. Then
\[
S := \sum_n\frac{\Psi(n)^2}{n}e^{-n/X} \le e^{-1/X} + D_4 + \mathrm{REM} + 2\sqrt{S_sS_t} + S_t,
\]
where
\[
D_4 = \Delta^{-2}\Big[\frac{Y_2^2 - Y_1^2}{2} - \gamma_E\Delta + k_1(Y_1+Y_2-2\gamma_E) + 2k_0\Big], \qquad k_1 \le -0.5853,\quad k_0 \le 0.4214 .
\]
Here $\mathrm{REM}$ is an explicit sum of exponentially small contour remainders, $S_t$ is an explicit bound for the mean square of the non-smooth part of the weights, and $S_s = e^{-1/X} + D_4 + \mathrm{REM}$.
\end{theorem}

\begin{proof}[Outline]
Write $\Psi = (\Lambda_V - \Lambda_U)/\Delta$, where $\Lambda_z(n) = \sum_{d\mid n}\mu(d)\log^+(z/d)$.
Split each truncation weight $\rho(d)$ as $\rho_s + \rho_t$, where $\rho_s(d) = \mathbb E[\rho(de^{-\eta Z})]$ with $Z\sim N(0,1)$ and $\eta = 0.02$.
By Minkowski's inequality in $\ell^2(e^{-n/X}/n)$, $S \le (\sqrt{S_s}+\sqrt{S_t})^2$.

For the smooth part, the Mellin transform of $\rho_s$ is $\hat\rho(s)e^{\eta^2s^2/2}$, and
\begin{multline*}
S_s = (2\pi i)^{-3}\iiint\Gamma(w)X^w\hat\rho_s(s_1)\hat\rho_s(s_2)\\ \times\frac{\zeta(1+w)\zeta(1+w+s_1+s_2)}{\zeta(1+w+s_1)\zeta(1+w+s_2)}\,\mathcal A(w,s_1,s_2)\,dw\,ds_1\,ds_2,
\end{multline*}
where $\mathcal A$ is an Euler product that converges absolutely near $0$ and satisfies $\mathcal A(w,0,s_2) = \mathcal A(w,s_1,0) = 1$.

The Gaussian factor allows every contour to be closed inside the box $|\Ip| \le T_2 = 400$, $\Rp \ge 1-\theta$, $\theta = 0.4$.
In this box $1/\zeta$ is analytic; this follows from \cite{PlattTrudgianRH} and is in any case confirmed by the computation.
Beyond the box all lines lie at $\Rp\ge1$, and the Gaussian factor gives $e^{-\eta^2T_2^2/2}$.

The residues produce $e^{-1/X} + D_4$, with $k_1 = H(0)$ and $k_0 = H'(0) + \gamma_EH(0)$, where $H$ is the kernel of the $s$-integral at $w = 0$.
At leading order this reproduces the asymptotic $\frac32 + (\Theta - \tau)/(\tau-1)$ of the Graham-type mean square.

The certified inputs are:
\begin{itemize}
\item rigorous box bounds $|\zeta| \le 11.89$ (for $|s-1|\ge 0.1$) and $|1/\zeta| \le 14.57$ on $\sigma\in[0.6,1.3]$, $|t|\le450$, obtained from Taylor models of $(s-1)\zeta(s)$ in \code{arb} (\code{lemG/box2.py});
\item enclosures $k_1 = -0.607063\pm0.0217$ and $k_0 = 0.35026\pm0.0711$ (\code{lemG/kcert2.py});
\item sup bounds for $\mathcal A$, obtained from enclosures over cells of the $3$-torus of phases for $p\le100$, a second-order bound for $p\le10^6$, and a tail estimate beyond (\code{lemG/remainder.py});
\item $S_t$, for which $|\rho_t(d)| \le (\eta/\Delta)\big[\varpi((\log d - \log U)/\eta) + \varpi((\log d - \log V)/\eta)\big]$ with $\varpi(y) = \phi(y) - |y|\bar\Phi(|y|)$.
  The resulting block sum uses exact masses for $m\le10^6$ and the bound $|\sum_{m\le x}\mu^2(m)/\varphi(m) - \log x - 1.3326| \le 4/\sqrt x$ \cite{RamareAkhilesh} beyond (\code{lemG/stbound.py}).
\end{itemize}
The integrals of explicit majorants are evaluated as rigorous upper sums or in closed form.
At the binding cell, $\mathrm{REM} = 5.4\cdot10^{-4}$ and $S_t = 9.9\cdot10^{-5}$.
All cells of the final tables have $\tau\in[1.54,1.78]$, $\log X/\log V\in[1.42,1.55]$ and $\log U\in[56.6,78.6]$, and every one of them takes its $d_4$ from Theorem~\ref{thm:G}.
Details are in \code{notes/LEMMA\_G.md}, Step~3.
\end{proof}

\begin{remark}
At the binding cell, Theorem~\ref{thm:G} gives $d_4 \approx 1.48$.
At the same parameters, an $X^\eps$ majorant combined with the Ramar\'e--Zuniga Alterman bound \cite{RZA} gives $\approx 20$, and our earlier explicit Selberg-diagonalisation bound (\code{notes/LEMMA\_F.md}) gives $\approx 5.8$; Bellotti's Table~5 values, at her parameters, are $85$--$105$.
Since $\log N$ has elasticity about $1$ in $d_4$, this is the single largest improvement in the density estimate.
\end{remark}
\section{The computation}\label{sec:computation}

\subsection{The scan}
For each interval $[L_1, L_2]$ we search for kernel parameters $(c, f)$:
first over the grid $c\in\{40,50,\dots,100\}$, $f\in\{0.5,\dots,0.9\}$, then locally over $c\pm10$ in steps of $5$ and $f\pm0.1$ in steps of $0.05$.
Given $(c,f)$, we put $\tau_{12} = \eps_R(L_2+1)/(1 + 2/L_1)$ and choose $\eps$ so that the smoothing term of Proposition~\ref{prop:master} equals $f\tau_{12}$.
This uses three fixed-point steps; the residual excess is bounded rigorously in \code{certify/transfer\_check.py} and is below $10^{-3}$ relative even at $L_1 = 43$.
We then evaluate the right-hand side of Proposition~\ref{prop:master} with $N^+$ from \S\ref{sec:density}, $\lambda^*$ from Theorem~\ref{thm:H} and $H = 3\cdot10^{12}$.
The ratio $q$ is this right-hand side divided by $\tau_{12}$; an interval is certified when $q < 1$.

The scan starts with intervals of width $10$ and bisects any interval that is not certified, down to width $0.5$ (\code{code/interval\_adapt.py}).
For $[43, 57)$ we only need some bound, since this range enters only through the integral in Lemma~\ref{lem:transfer}; there we use widths $2$ to $0.5$, and $q$ may exceed $1$.

\begin{center}
\begin{tabular}{lll}
\toprule
range of $L$ & worst $q$ & remark\\
\midrule
$[43,57)$ & $3.80$ (at $43$), $\le 0.47$ above $48$ & used only through Lemma~\ref{lem:transfer}\\
$[57,1000)$ & $0.9966$ at $907$ & $[857,907)$ needed width $5$\\
$[1000,2000)$ & $0.8884$ at $1000$ & \\
$[2000,2810)$ & $0.4063$ & \\
$[2800,3159)$ & $0.4001$ & \\
\bottomrule
\end{tabular}
\end{center}

The binding range $L\approx 900$ corresponds to bandwidth $T_1 \approx e^{30}$.
There the zero sum is dominated by $\lambda = (1-\beta)\log T \in [0.20, 0.32]$ and $\log T\in(28.73, 30]$.
For the interval $[900,910]$ of width $10$, the best kernel gives $0.40$ (smoothing) $+ 0.60$ (zeros above $H_0$) $= 1.002$, plus negligible terms; intervals of width $5$ are certified.

\subsection{Certificates}
The following list gives, for each ingredient, the program that certifies it.
All these programs are in the repository; \code{code/rerun\_all.sh} re-runs the whole chain.
\begin{enumerate}
\item Lemma~\ref{lem:epsR} ($\eps_R$ on every cell): \code{certify/epsR\_cert.py}.
\item Lemma~\ref{lem:smooth} (smoothing constant): \code{certify/smoothing\_const.py}.
\item Proposition~\ref{prop:transfer} (from $\psi$ to $\pi - \li$, both conditions of \eqref{eq:target}): \code{certify/transfer\_check.py}.
  The worst ratio is $0.99577$, at $\log y = 917$.
\item Lemmas~\ref{lem:W0}--\ref{lem:L12} (constants of \cite{BTY}): in \code{rep/}, the programs \code{consts\_cert.py}, \code{lemma5\_cert.py}, \code{lemma11\_cert.py}, \code{f\_nonneg\_cert.py}, \code{mp\_cert.py} and \code{misc\_cert.py}.
\item Theorem~\ref{thm:H}: \code{rep/lamcert.py}, with output \code{rep/lamstar2c.json}.
\item Theorem~\ref{thm:rep}: \code{rep/repcert\_run.py} and \code{rep/min\_fix.py}; output \code{rep/dstar2c\_all.json}.
\item Theorem~\ref{thm:G}: in \code{lemG/}, the programs \code{box2.py}, \code{kcert2.py}, \code{remainder.py} and \code{stbound.py}.
\item Theorem~\ref{thm:density}: the tables \code{lemmaD\_tableGH0.01\_*.pkl}, produced by \code{bt4.py}.
  Every cell's bound is recomputed exactly from its stored parameters when it is used.
\item The scan: \code{interval\_adapt.py}, with logs in \code{logs/}.
\end{enumerate}

\subsection{Rounding in the scan}
Items 1--7 are certified in ball arithmetic.
The scan (item~9) and the evaluation of the density bounds from the stored parameters (item~8) run in IEEE double precision, with $\eps_R$ and $-\zeta'/\zeta(1+\eta)$ taken from $60$- and $30$-digit \code{mpmath} computations.
This is the one place where the argument rests on a floating-point computation rather than on a certificate, so we describe it precisely.

Each value of $q$ is a sum of a few thousand non-negative terms.
Each term is a product of exponentials and explicit factors, obtained from the parameters by a bounded number (well under $10^4$) of floating-point operations.
Only two subtractions involve quantities of comparable size.
The first is $c_1 = e^{-L^2/U} - B_1 - \cdots$; over all table cells $c_1 \ge 0.92$, so there is no cancellation.
The second is $c_1^2 - c_5$, which is $\ge 0.83$ over all table cells, so it amplifies relative errors by at most a factor of about $1.2$.
The greedy knapsack value of Lemma~\ref{lem:pair} and the Stieltjes and upper sums of Proposition~\ref{prop:master} are sums of non-negative terms.
Standard error analysis therefore bounds the relative rounding error of each $q$ by about $10^{-11}$.

As an empirical check, inflating every density bound by a factor $1+10^{-6}$ changes $q$ on the binding interval $[907,912]$ by $5.8\cdot10^{-7}$.
This is consistent with the scan depending on the density bounds linearly, with coefficient at most $1$.
The smallest margins are $1-q \ge 3\cdot10^{-3}$ in the scan and about $4\cdot10^{-3}$ in Proposition~\ref{prop:transfer}, some eight orders of magnitude above the rounding error.
A complete interval-arithmetic re-evaluation of the scan, keeping the floating-point parameters but evaluating each bound in \code{arb}, is straightforward but costly; we have not carried it out.

\subsection{Adversarial audits}
Two adversarial audits were carried out by separate instances of the same AI system that assisted with the computations (see the acknowledgements); they are not independent in the human sense.
The first audit covered the new ingredients (\S\S\ref{sec:zfr}--\ref{sec:density}) and checked the following points, confirming them:
\begin{itemize}
\item the legitimacy of keeping the second pair in Lemma~6 of \cite{BTY};
\item the per-$\sigma$ form of Lemma~5;
\item the induction in Theorem~\ref{thm:H};
\item the use of $d^*$ in the pair sum;
\item the error constants;
\item the replacement of $\mathcal J$.
\end{itemize}
It found one genuine defect: the clamp on $\nu$ described at the end of \S\ref{sec:rep}.
It also found three minor gaps: a missing non-negativity certificate for the dropped Lemma-10 terms, the unrounded $w$-derivative values, and the range of the Hiary--Patel--Yang bound.
All four were fixed, and the full chain was re-run with identical results.

A second audit, of a draft of this paper, found an unjustified rescaling: the density constants, which are valid at the corner $\delta_k$ of a cell, had been applied with the actual $\delta$ inside the cell.
This was removed (Theorem~\ref{thm:density} now uses the corner value), and the scan was re-run.
That audit also corrected the definition of $N_R$, the attributions for the lower range, the hypotheses of Lemmas~\ref{lem:L5} and~\ref{lem:epsR}, the statement of the correction factor $\iota$, and the reproducibility scripts.
\section{Remarks}\label{sec:remarks}

\subsection{Corrections to inputs}
For convenience we collect the points found while re-deriving the explicit inputs; none of them changes a published main theorem.
\begin{enumerate}
\item In \cite{BTY}:
  \begin{enumerate}
  \item the Lemma-10 term in the final inequality should carry the factor $\kappa$ (Remark~\ref{rem:btyf0});
  \item the misprint in (25) (Remark~\ref{rem:btyf0});
  \item the value of $m_5$ in Table~1 (Remark~\ref{rem:m5});
  \item the integral of the digamma majorant (31) (\S\ref{ss:constants});
  \item the uniform test for (21), which needs $\sigma \ge 0.993464$ (Remark~\ref{rem:L5});
  \item the numerical verification of $\partial_\sigma G_p < 0$ is unnecessary, since this follows from $P \ge 0$ (Lemma~\ref{lem:L1011});
  \item the terms with $m > 15$ in Lemma~10 are dropped there without a bound, which is justified by $f\ge0$.
  \end{enumerate}
\item In B\"uthe's smoothing estimate, the quoted asymptotic for $\nu_c$ does not match the stated normalisation (Remark~\ref{rem:buthe}).
\item In Bellotti \cite{Bellotti}, the pair-sum accounting \cite[(4.11)]{Bellotti} assumes a spacing of $\delta$ between representatives of different rectangles, which does not appear to be justified as stated; Lemma~\ref{lem:pair} avoids this assumption.
  Since her constants are much larger than ours elsewhere, her published bounds are not affected in practice.
\end{enumerate}

\subsection{On two recent preprints}\label{sec:de}
\begin{enumerate}
\item \emph{arXiv:2408.02591} \cite{De2} claims \eqref{eq:ram} for $x \ge e^{43.5102147}$.
  The argument combines $\pi(x) < e\,\pi(x/e)$ with $\pi(x) > x/\log x$ to conclude, in (29), that $\pi(x)^2 - \frac{ex}{\log x}\pi(x/e) < 0$.
  That conclusion requires $\frac{x}{\log x}e\pi(x/e) \ge \pi(x)^2$.
  From the first inequality one only obtains $\frac{x}{\log x}e\pi(x/e) > \frac{x}{\log x}\pi(x)$.
  Passing from there to $\pi(x)^2$ would need $\pi(x) \le x/\log x$, which is the reverse of the second inequality (and is false for $x \ge 17$).
  The gap is not a slip that can be repaired: the difference of the two sides of \eqref{eq:ram} has size $x^2/\log^6x$, whereas the bounds used have relative accuracy $\approx(\log x)^{-3/2}$.
\item \emph{arXiv:2407.12052} \cite{De1} claims \eqref{eq:ram} for $x\ge e^{547}$.
  The argument has four problems.
  \begin{enumerate}
  \item The monotonicity argument assumes Sterneck's conjecture $|M(x)| < \frac12\sqrt x$, which is false \cite{Neubauer}.
  \item It asserts $M(x) = O(\sqrt x)$, which is an open problem (and widely believed false); the proof given yields at most $O(x^{1/2+\eps})$.
  \item The retained main term is itself only $O(x^2/\log^2x)$, so the sign of $G(x) = \pi(x)^2 - \frac{ex}{\log x}\pi(x/e)$ is not determined.
  \item The range below $e^{3158.442}$ is "covered" by evaluating at $26$ sample points.
  \end{enumerate}
\end{enumerate}
Neither preprint therefore proves \eqref{eq:ram} in the range considered here.

\subsection*{Acknowledgements}
The computations and the drafting of this paper were carried out with substantial assistance from an AI system (Claude, Anthropic), which also performed the adversarial audit described in \S\ref{sec:computation}.
All certificates are reproducible from the public code.
\begin{thebibliography}{99}
\bibitem{arb} F.~Johansson, \emph{Arb: efficient arbitrary-precision midpoint-radius interval arithmetic}, IEEE Trans. Comput. \textbf{66} (2017), 1281--1292.
\bibitem{Axler2} C.~Axler, \emph{Estimates for $\pi(x)$ for large values of $x$ and Ramanujan's prime counting inequality}, Integers \textbf{18} (2018), A61.
\bibitem{Axler} C.~Axler, \emph{On Ramanujan's prime counting inequality}, arXiv:2207.02486.
\bibitem{Bellotti} C.~Bellotti, \emph{An explicit log-free zero density estimate for the Riemann zeta-function}, arXiv:2405.12545.
\bibitem{BTY} C.~Bellotti, T.~Trudgian and A.~Yang, \emph{Zero-free regions inspired by work of Heath-Brown}, arXiv:2603.21490.
\bibitem{BerndtIV} B.~C.~Berndt, \emph{Ramanujan's Notebooks, Part IV}, Springer, New York, 1994.
\bibitem{Buthe2016} J.~B\"uthe, \emph{Estimating $\pi(x)$ and related functions under partial RH assumptions}, Math. Comp. \textbf{85} (2016), 2483--2498.
\bibitem{Buthe2018} J.~B\"uthe, \emph{An analytic method for bounding $\psi(x)$}, Math. Comp. \textbf{87} (2018), 1991--2009.
\bibitem{CullyHugillJohnston} M.~Cully-Hugill and D.~R.~Johnston, \emph{On the error term in the explicit formula of Riemann--von Mangoldt}, Int. J. Number Theory \textbf{19} (2023), 1205--1228.
\bibitem{De1} S.~De, \emph{On proving an inequality of Ramanujan using explicit order estimates for the Mertens function}, arXiv:2407.12052.
\bibitem{De2} S.~De, \emph{On proving Ramanujan's inequality using a sharper bound for the prime counting function $\pi(x)$}, arXiv:2408.02591.
\bibitem{DudekPlatt} A.~W.~Dudek and D.~J.~Platt, \emph{On solving a curious inequality of Ramanujan}, Exp. Math. \textbf{24} (2015), 289--294.
\bibitem{Hassani} M.~Hassani, \emph{On an inequality of Ramanujan concerning the prime counting function}, Ramanujan J. \textbf{28} (2012), 435--442.
\bibitem{HSW} E.~Hasanalizade, Q.~Shen and P.-J.~Wong, \emph{Counting zeros of the Riemann zeta function}, J. Number Theory \textbf{235} (2022), 219--241.
\bibitem{Hiary} G.~A.~Hiary, \emph{An explicit van der Corput estimate for $\zeta(1/2+it)$}, Indag. Math. \textbf{27} (2016), 524--533.
\bibitem{HPY} G.~A.~Hiary, D.~Patel and A.~Yang, \emph{An improved explicit estimate for $\zeta(1/2+it)$}, J. Number Theory \textbf{256} (2024), 195--217.
\bibitem{Johnston} D.~R.~Johnston, \emph{Improving bounds on prime counting functions by partial verification of the Riemann hypothesis}, Ramanujan J. \textbf{59} (2022), 1307--1321; arXiv:2109.02249.
\bibitem{JohnstonYang} D.~R.~Johnston and A.~Yang, \emph{Some explicit estimates for the error term in the prime number theorem}, J. Math. Anal. Appl. \textbf{527} (2023), 127460.
\bibitem{KLN} H.~Kadiri, A.~Lumley and N.~Ng, \emph{Explicit zero density for the Riemann zeta function}, J. Math. Anal. Appl. \textbf{465} (2018), 22--46.
\bibitem{MV73} H.~L.~Montgomery and R.~C.~Vaughan, \emph{The large sieve}, Mathematika \textbf{20} (1973), 119--134.
\bibitem{Neubauer} G.~Neubauer, \emph{Eine empirische Untersuchung zur Mertensschen Funktion}, Numer. Math. \textbf{5} (1963), 1--13.
\bibitem{PY} D.~Patel and A.~Yang, \emph{An explicit sub-Weyl bound for $\zeta(1/2+it)$}, J. Number Theory (2024).
\bibitem{PlattTrudgian} D.~J.~Platt and T.~S.~Trudgian, \emph{The error term in the prime number theorem}, Math. Comp. \textbf{90} (2021), 871--881.
\bibitem{PlattTrudgianRH} D.~J.~Platt and T.~S.~Trudgian, \emph{The Riemann hypothesis is true up to $3\cdot10^{12}$}, Bull. Lond. Math. Soc. \textbf{53} (2021), 792--797.
\bibitem{RamareAkhilesh} O.~Ramar\'e and P.~Akhilesh, \emph{Explicit averages of non-negative multiplicative functions: going beyond the main term}, Colloq. Math. \textbf{147} (2017), 275--313.
\bibitem{RZA} O.~Ramar\'e and S.~Zuniga Alterman, \emph{An $L^2$-bound for the Barban--Vehov weights}, arXiv:2405.12662.
\bibitem{RS62} J.~B.~Rosser and L.~Schoenfeld, \emph{Approximate formulas for some functions of prime numbers}, Illinois J. Math. \textbf{6} (1962), 64--94.
\end{thebibliography}
\end{document}
